\section{Medición del Factor de Calidad}

\subsection{Cálculo Analítico}

El factor de calidad $Q$ para un circuito RLC serie se define como:

\begin{equation}
    Q = \frac{1}{R}\sqrt{\frac{L}{C}} = \frac{\omega_0 L}{R}
    \label{eq:Q_def}
\end{equation}

Utilizando los valores nominales de los componentes 
($L = \SI{38.76}{\micro\henry}$, $C = \SI{8.2}{\nano\farad}$, $R = \SI{68}{\ohm}$):

\begin{equation}
    Q_{\text{nominal}} = \frac{1}{68}\sqrt{\frac{38{,}76 \times 10^{-6}}{8{,}2 \times 10^{-9}}} = 1{,}00
    \label{eq:Q_nominal}
\end{equation}

Utilizando los valores medidos experimentalmente 
($L = \SI{38.3}{\micro\henry}$ a $f = \SI{200}{\kilo\hertz}$, $C = \SI{8.58}{\nano\farad}$):

\begin{equation}
    Q_{\text{teórico}} = \frac{1}{68}\sqrt{\frac{38{,}3 \times 10^{-6}}{8{,}58 \times 10^{-9}}} = 1{,}022
    \label{eq:Q_teo}
\end{equation}

\subsection{Métodos Empíricos y Selección del de Menor Error}

Se analizaron dos métodos para medir $Q$ experimentalmente con el osciloscopio 
y el generador de señales.

\subsubsection*{Método A — Frecuencia del pico de amplitud}

En la configuración pasabajos (salida en el capacitor), la ganancia presenta un 
máximo no en $f_0$ sino en una frecuencia ligeramente inferior, dada por:

\begin{equation}
    f_{\text{pico}} = f_0 \sqrt{1 - \frac{1}{2Q^2}}
    \label{eq:f_pico}
\end{equation}

Midiendo experimentalmente $f_{\text{pico}}$ y conociendo $f_0$, se puede 
despejar $Q$ de la ecuación~\ref{eq:f_pico}. La ventaja de este método es que 
$f_{\text{pico}}$ es localizable ajustando la frecuencia del generador hasta 
maximizar la amplitud de $V_{out}$. Sin embargo, cualquier imprecisión en la 
determinación de $f_0$ se propaga directamente al cálculo de $Q$, siendo este 
su principal fuente de error.

\subsubsection*{Método B — Razón de amplitudes en resonancia}

En la frecuencia de resonancia $f_0$, la función de transferencia del capacitor 
tiene módulo igual a $Q$:

\begin{equation}
    \left|H(j\omega_0)\right| = Q \implies Q = \frac{\left|V_{out}(f_0)\right|}{\left|V_{in}(f_0)\right|}
    \label{eq:Q_metodo_B}
\end{equation}

Este método requiere únicamente medir las amplitudes de $V_{in}$ y $V_{out}$ 
simultáneamente en $f_0$ con los dos canales del osciloscopio. Al tratarse de 
un cociente de amplitudes en un único punto de frecuencia, el error queda 
limitado a la resolución vertical del instrumento, sin depender de un barrido 
fino de frecuencias. Por lo anterior, el Método B es el que introduce 
menor error y fue adoptado como método principal.

Para localizar $f_0$ con mayor precisión previo a aplicar el Método B, se 
utilizó la condición de fase: en $\omega_0$ la tensión en el capacitor se 
encuentra desfasada exactamente $-90°$ respecto de $V_{in}$. Alimentando el 
circuito con una señal sinusoidal de \SI{1}{V_{pp}}, se varió la frecuencia 
hasta observar en el osciloscopio que ambas sinusoides presentaban un desfasaje 
de aproximadamente $90°$, como se muestra en la figura~\ref{fig:desfasaje_f0}.

\subsection{Resultados Experimentales}

\subsubsection*{Determinación de \texorpdfstring{$f_0$}{f0}}

Mediante el procedimiento de desfasaje descripto se obtuvo:

\begin{equation}
    f_{0,\,\text{medido}} = \SI{273}{\kilo\hertz} \qquad 
    f_{0,\,\text{teórico}} = \frac{1}{2\pi\sqrt{LC}} = \SI{282.7}{\kilo\hertz}
    \label{eq:f0_comparacion}
\end{equation}

La diferencia del $3{,}4\%$ es consistente con las tolerancias de los 
componentes, en particular el capacitor de poliéster con tolerancia del $20\%$.

\begin{figure}[H]
    \centering
    %\includegraphics[width=0.7\textwidth]{imagenes/desfasaje_f0.png}
    \caption{Desfasaje de $\approx 90°$ entre $V_{in}$ (amarillo) y $V_{out}$ 
             (verde) a $f_0 = \SI{273}{\kilo\hertz}$, confirmando la 
             frecuencia de resonancia.}
    \label{fig:desfasaje_f0}
\end{figure}

\subsubsection*{Método A}

Buscando el máximo de amplitud de $V_{out}$ mediante barrido de frecuencia, 
se halló la frecuencia del pico:

\begin{equation}
    f_{\text{pico}} \approx \SI{197}{\kilo\hertz}
    \label{eq:f_pico_medido}
\end{equation}

Despejando $Q$ de la ecuación~\ref{eq:f_pico} con los valores de las 
ecuaciones~\ref{eq:f0_comparacion} y~\ref{eq:f_pico_medido}:

\begin{equation}
    Q_A = \frac{1}{\sqrt{2\left[1-\left(\dfrac{f_{\text{pico}}}{f_0}\right)^2\right]}} 
        = \frac{1}{\sqrt{2\left[1-\left(\dfrac{197}{273}\right)^2\right]}} 
        = 0{,}995
    \label{eq:QA}
\end{equation}

\begin{figure}[H]
    \centering
    %\includegraphics[width=0.7\textwidth]{imagenes/metodo_A_pico.png}
    \caption{Medición de $f_{\text{pico}} \approx \SI{197}{\kilo\hertz}$ 
             en el osciloscopio. Se observa la amplitud máxima de $V_{out}$ 
             (canal 1, verde) respecto de $V_{in}$ (canal 2, amarillo).}
    \label{fig:metodo_A}
\end{figure}

\subsubsection*{Método B}

Se midieron las amplitudes pico a pico de $V_{out}$ y $V_{in}$ en $f_0$ 
con los cursores del osciloscopio, tal como se ilustra en la 
figura~\ref{fig:metodo_B}. Aplicando la ecuación~\ref{eq:Q_metodo_B}:

\begin{equation}
    Q_B = \frac{\left|V_{out}(f_0)\right|}{\left|V_{in}(f_0)\right|} = 1{,}05
    \label{eq:QB}
\end{equation}

\begin{figure}[H]
    \centering
    %\includegraphics[width=0.7\textwidth]{imagenes/metodo_B_amplitudes.png}
    \caption{Medición de amplitudes en $f_0 = \SI{273}{\kilo\hertz}$ 
             mediante cursores. Canal 1 (verde): $V_{out}$; 
             Canal 2 (amarillo): $V_{in}$.
             El cociente de amplitudes da $Q_B = 1{,}05$.}
    \label{fig:metodo_B}
\end{figure}

\subsection{Comparación y Análisis}

En la tabla~\ref{tab:Q_comparacion} se resumen los valores obtenidos 
por cada método.

\begin{table}[H]
    \centering
    \caption{Comparación de valores del factor de calidad $Q$.}
    \label{tab:Q_comparacion}
    \begin{tabular}{lc}
        \toprule
        Método & $Q$ \\
        \midrule
        Nominal (valores de catálogo)       & $1{,}00$ \\
        Teórico (componentes medidos)       & $1{,}022$ \\
        Método A (frecuencia del pico)      & $0{,}995$ \\
        Método B (razón de amplitudes)      & $1{,}05$  \\
        \bottomrule
    \end{tabular}
\end{table}

Ambos métodos arrojan resultados dentro del $4\%$ del valor teórico calculado 
con los componentes medidos, lo cual es aceptable ya que el inductor presenta una resistencia 
serie parásita $r_L$ que incrementa el $R$ efectivo del circuito, reduciendo el 
$Q$ real respecto del nominal.


El método B es el más práctico y repetible debido a que introduce menor error. En particular, en este
escenario solo basta con hacer el cociente de las amplitudes en el entorno cercano a la fracuencia natural
del sistema. Por otra parte, el método A induce mayor error debido a que es más sensible respecto a la 
variable del $f?0$. Esto se denota en la expresión~\ref{eq:f_pico}, la cual involucra una diferencia de
cuadrados de frecuencias relativamente cercanas.

